% !TEX program = pdflatex
\documentclass[11pt,a4paper]{article}

% =========================================================================
% ACTA DE EXCLUSIÓN DE HORUS SENTINEL
% Documento común del equipo
% Gestión de Proyectos TICS (EST315) -- Semestre 2026-II -- Team Sobrevivientes
% =========================================================================

% --- PAQUETES DE CONFIGURACIÓN ---
\usepackage[utf8]{inputenc}
\usepackage[T1]{fontenc}
\usepackage[spanish,es-tabla]{babel}
% Compatibilidad con babel-spanish reciente (flechas de TikZ y signo de porcentaje)
\AtBeginDocument{\shorthandoff{<>}\spanishplainpercent}
\usepackage{geometry}
\geometry{
    top=3.6cm,
    bottom=2.3cm,
    left=2.2cm,
    right=2.2cm,
    % PARTE MODIFICABLE: headheight debe ser mayor que la altura del logo más alto (2.56cm = 73pt).
    headheight=80pt,
    headsep=15pt,
    footskip=20pt
}

\usepackage{graphicx}
\usepackage{fancyhdr}
\usepackage[table]{xcolor}
\usepackage{titlesec}
\usepackage{booktabs}
\usepackage{tabularx}
\usepackage{array}
\usepackage{amsmath}
\usepackage{enumitem}
\usepackage{microtype}
\usepackage{amssymb}
\usepackage{needspace}
\usepackage{hyperref}

% Logos: funciona a 2 o 3 niveles de profundidad dentro de docs/latex/
\graphicspath{{../../assets/img/}{../../../assets/img/}}

% --- PALETA INSTITUCIONAL FIEI / UNA PUNO ---
\definecolor{AzulUNA}{RGB}{10, 45, 90}
\definecolor{DoradoFIEI}{RGB}{180, 130, 20}
\definecolor{GrisOscuro}{RGB}{50, 50, 50}
\definecolor{AzulClaro}{RGB}{228, 236, 246}

% --- CONFIGURACIÓN DE LISTAS ---
\setlist{topsep=2pt, itemsep=2pt, parsep=1pt}

% --- COLUMNAS DE TABLA ---
\newcolumntype{L}{>{\raggedright\arraybackslash}X}
\newcolumntype{P}[1]{>{\raggedright\arraybackslash}p{#1}}
\newcolumntype{C}[1]{>{\centering\arraybackslash}p{#1}}
\newcolumntype{Q}[1]{>{\raggedleft\arraybackslash}p{#1}}
\newcolumntype{R}{>{\raggedleft\arraybackslash}X}
\renewcommand{\arraystretch}{1.25}

% --- FORMATO DE ENLACES ---
\hypersetup{
    colorlinks=true,
    linkcolor=AzulUNA,
    urlcolor=AzulUNA,
    citecolor=AzulUNA,
    pdftitle={Acta de Exclusión de Horus Sentinel},
    pdfauthor={Team Sobrevivientes}
}

% --- FORMATO DE TÍTULOS ---
\titleformat{\section}{\color{AzulUNA}\normalfont\large\bfseries}{\thesection.}{0.7em}{}[\color{AzulUNA}\titlerule]
\titleformat{\subsection}{\color{AzulUNA}\normalfont\normalsize\bfseries}{\thesubsection.}{0.5em}{}
\titlespacing*{\section}{0pt}{9pt}{4pt}
\titlespacing*{\subsection}{0pt}{6pt}{2pt}

% --- ENCABEZADO Y PIE DE PÁGINA (fancyhdr) ---
\pagestyle{fancy}
\fancyhf{}

% ===== TAMAÑO DE LOS LOGOS (PARTE MODIFICABLE) =====
% logo_fiei.png  : height=1.45cm (lado izquierdo)
% logo_equipo.png: height=2.56cm (lado derecho; 25 % más grande porque la imagen tiene más borde blanco)
% Para cambiar el tamaño, modifique solo el valor height=... de cada \includegraphics.
% Si el logo más alto supera 2.6cm, aumente también headheight y top en la configuración de página.
\fancyhead[L]{%
    \includegraphics[height=1.45cm,keepaspectratio]{logo_fiei.png}%
}
\fancyhead[C]{%
    \scriptsize \textbf{\textcolor{AzulUNA}{UNIVERSIDAD NACIONAL DEL ALTIPLANO DE PUNO}}\\[2pt]
    \tiny \textcolor{GrisOscuro}{FACULTAD DE INGENIERÍA ESTADÍSTICA E INFORMÁTICA}\\[2pt]
    \scriptsize \textbf{\textcolor{AzulUNA}{Gestión de Proyectos TICS (EST315) \textbar\ Team Sobrevivientes}}\\[2pt]
    \footnotesize \textbf{\textcolor{DoradoFIEI}{ACTA DE EXCLUSIÓN DE HORUS SENTINEL}}%
}
\fancyhead[R]{%
    \includegraphics[height=2.56cm,keepaspectratio]{logo_equipo.png}%
}

\renewcommand{\headrulewidth}{0pt}
\renewcommand{\headrule}{}

\fancyfoot[L]{\scriptsize \textcolor{GrisOscuro}{CasaScore \textperiodcentered\ Horus Sentinel}}
\fancyfoot[C]{\small \thepage}
\fancyfoot[R]{\scriptsize \textcolor{GrisOscuro}{Puno, Perú -- Semestre 2026-II}}
\renewcommand{\footrulewidth}{0.4pt}

\begin{document}

\begin{center}
    {\Large \textbf{\textcolor{AzulUNA}{ACTA DE EXCLUSIÓN DE HORUS SENTINEL}}}\\[3pt]
    {\normalsize \textbf{Exclusión del proyecto, liberación de horas y cierre de su expediente}}\\[4pt]
    \footnotesize \textbf{Fecha:} martes 22/09/2026 (semana 5) \quad \textbar\quad \textbf{Hora:} 20:00 -- 21:00 \quad \textbar\quad \textbf{Decisión:} unanimidad del equipo
\end{center}

\vspace{0.1cm}

\noindent\begin{tabularx}{\textwidth}{@{}P{3.8cm} L@{}}
\toprule
\textbf{Curso:} & Gestión de Proyectos TICS (EST315) -- Semestre 2026-II \\
\textbf{Equipo:} & Team Sobrevivientes, sponsor de ambos proyectos (los cuatro integrantes presentes) \\
\textbf{Proyecto excluido:} & Horus Sentinel \\
\textbf{Proyecto que continúa:} & CasaScore (AlquilaPuno) \\
\textbf{Insumos revisados:} & Informe de control de Horus Sentinel al corte de la semana 4, comparación financiera de la cartera, Business Case, cronograma y presupuesto de Horus Sentinel \\
\bottomrule
\end{tabularx}

\section{Antecedentes}
Desde la semana 3, el equipo gestiona dos proyectos en paralelo con 8 horas semanales por integrante en cada uno. El Acta de Evaluación de Cartera dio prioridad a CasaScore y aprobó Horus Sentinel con alcance ajustado. El Acta de Constitución de Horus fijó como riesgos de mayor severidad el tiempo y el presupuesto, a revisar con el cronograma y el presupuesto detallados de la semana 4. Al aprobarse ambos, el Project Manager encargó medir el estado real de Horus Sentinel y comparar su rentabilidad con la de CasaScore. Los resultados se presentaron en esta reunión.

\section{Motivos de la Exclusión}
{\small
\noindent\begin{tabularx}{\textwidth}{@{}P{3.0cm} L@{}}
\toprule
\rowcolor{AzulClaro}
\textbf{Motivo} & \textbf{Evidencia} \\
\midrule
\textbf{Tiempo insuficiente} & El cronograma detallado de Horus requiere 19 semanas de trabajo con 8 horas semanales por integrante; la probabilidad de terminar en la semana 17 es menor a 0.1\,\%. Entre las semanas 5 y 17 quedan 13 semanas. \\
\addlinespace[2pt]
\textbf{Desempeño real} & Al corte de la semana 4: SPI de 0.79 y CPI de 0.94. La capacitación geoespacial va al 55\,\% y retrasa el modelo de datos de la ruta crítica. Al ritmo actual, el proyecto terminaría el 12/02/2027 con un costo de hasta S/ 15,406.80. \\
\addlinespace[2pt]
\textbf{Presupuesto alto} & El presupuesto detallado de Horus asciende a S/ 13,663.57, con S/ 3,229.30 de gasto en efectivo. Sumado al de CasaScore, cada integrante debe aportar S/ 1,434.19 en el semestre (S/ 807.33 de Horus y S/ 626.86 de CasaScore). \\
\addlinespace[2pt]
\textbf{Carga académica} & Mantener ambos proyectos exige 16 horas semanales por integrante, además de las demás asignaturas del semestre. En las semanas 2 a 4 se registraron 97 de las 115 horas planificadas para Horus. \\
\addlinespace[2pt]
\textbf{Menor rentabilidad} & VAN económico de S/ +17,557.81 frente a S/ +64,059.98 de CasaScore, con una relación beneficio/costo de 2.48 frente a 6.77. Sin el convenio municipal, que aún no se ha gestionado, el VAN económico de Horus baja a S/ +2,636.71. \\
\bottomrule
\end{tabularx}
}

\Needspace{16\baselineskip}
\section{Estado de Horus Sentinel al Cierre}
{\small
\noindent\begin{tabularx}{\textwidth}{@{}C{2.2cm} L C{2.4cm}@{}}
\toprule
\rowcolor{AzulClaro}
\textbf{N.\textsuperscript{o}} & \textbf{Entregable} & \textbf{Estado} \\
\midrule
01 & Ficha técnica de la propuesta & Completo \\
02 & Especificación de requerimientos & Completo \\
03 & Viabilidad y Business Case & Completo \\
04 & Acta de Constitución & Completo \\
05 & Registro de interesados & Completo \\
06 & EDT y diccionario de la EDT & Completo \\
07 & Plan del equipo & Completo \\
08 & Cronograma & Completo \\
09 & Presupuesto & Completo \\
16 & Control del proyecto al corte de la semana 4 & Completo \\
10--15 y 17--28 & Demás entregables de planificación, ejecución y cierre & No se elaboran \\
\bottomrule
\end{tabularx}
}

\section{Resoluciones}
\begin{enumerate}[leftmargin=*]
    \item \textbf{Exclusión:} se aprueba la alternativa recomendada por el Project Manager y se excluye Horus Sentinel de la cartera del equipo a partir del 22/09/2026. Se descartaron continuar con el alcance completo, porque termina fuera del semestre, y reducirlo a un prototipo con datos simulados, porque sin reportes ciudadanos el proyecto no genera beneficio.
    \item \textbf{Archivo:} los entregables 01 a 09 y el informe de control quedan en la carpeta \texttt{horus\_sentinel/} del repositorio como evidencia del trabajo realizado. El proyecto no se retoma.
    \item \textbf{Cierre administrativo:} se cierran los Issues pendientes de Horus, se archiva su rama de trabajo y se crea la etiqueta \texttt{horus-archivo-2026}.
    \item \textbf{Horas liberadas:} desde la semana 6, cada integrante mantiene 8 horas semanales en CasaScore. Las 8 horas que dedicaba a Horus Sentinel quedan a libre disposición de cada integrante.
    \item \textbf{Aporte económico:} se anula el aporte destinado a Horus (S/ 807.33 por integrante); se mantiene el aporte para CasaScore.
    \item \textbf{CasaScore:} el cronograma y el presupuesto de CasaScore se mantienen; su Plan de Dirección registra esta exclusión.
    \item \textbf{Comunicación:} la presente acta se informa al docente asesor y se publica en el repositorio.
\end{enumerate}

\Needspace{20\baselineskip}
\section{Firmas}
\noindent En señal de conformidad con las resoluciones adoptadas por unanimidad, los integrantes del equipo firman la presente acta:

\vspace{1.5cm}
\begin{center}
\begin{tabular}{c c c}
    \parbox{6.5cm}{\centering \rule{6.0cm}{0.6pt}\\[5pt] \textbf{Clyde Neil Paricahua Pari}\\[2pt] {\footnotesize Project Manager}}
    & \hspace{0.8cm} &
    \parbox{6.5cm}{\centering \rule{6.0cm}{0.6pt}\\[5pt] \textbf{Ruth Karina Apaza Solis}\\[2pt] {\footnotesize Analista de Planificación}}
    \\[2.0cm]
    \parbox{6.5cm}{\centering \rule{6.0cm}{0.6pt}\\[5pt] \textbf{Eliseo Tarqui Ajahuana}\\[2pt] {\footnotesize Especialista Técnico / Arquitectura}}
    & \hspace{0.8cm} &
    \parbox{6.5cm}{\centering \rule{6.0cm}{0.6pt}\\[5pt] \textbf{Viviana Mary Quisbert Quispe}\\[2pt] {\footnotesize Gestora de Comunicación, Riesgos y Documentación}}
\end{tabular}
\end{center}
\end{document}
